Esta etapa se reparte en los siguientes archivos:

\begin{itemize}
\item \texttt{app/services/prediccion/modelos.py}
\item \texttt{app/services/prediccion/regresion\_lineal.py}
\end{itemize}

\subparagraph{Definicion de variables y funciones basicas}

Al incio del archivo \texttt{modelos.py}, se define el nivel de significancia \texttt{NIVEL\_SIGNIFICANCIA\_DM} como 5\%, esta variable se va a usar posteriormente a la hora de evaluar si un resultado es estadisticamente significativo:

\begin{minted}{python}
NIVEL_SIGNIFICANCIA_DM = 0.05
\end{minted}

Como parte de la propuesta de Thary,tambien se implemenaron las siguientes cuatro variables:

\begin{minted}{python}
MESES_ARCHIVO_UMBRAL_TEST_GRANDE = 18
TEST_SIZE_MESES_DEFAULT = 3
TEST_SIZE_MESES_GRANDE = 6

MESES_PERDIDOS_POR_REZAGOS = 3
\end{minted}

Por defecto, el periodo de prueba son 3 meses, pero cuando se cuenta con un periodo de prueba mas grande, de 18 meses o mas, este pasa de 3 meses a 6 meses.

Esta decision se tomo despues de identififcar un problema con el que contaba el sistema. Como punto de partida, el sistema elige el mejor modelo para cada producto comparando el MAE de los 3 modelos sobre 3 meses del historial, el problema era que con solo 3 meses, el resultado podria no ser lo suficientemente confiable, esto debido a que con tan pocos datos, un solo mes que sea atipico puede terminar haciendo que un modelo "gane" sin que eso refleje una diferencia real, sobretodo porque el MAE es una medicion que cuenta con mucha varianza. Entre los riesgos identificados se encuentra: poca capacidad de distinguir modelos parecidos, inestabilidad del resultado y sesgo estacional.

Antes de llegar a la solucion definitva se planteo una primera propuesta, en la que si el archivo contaba con menos de 18 meses, se haria la comparacion del MAE de los 3 modelos sobre 3 meses, pero si contaba con 18 meses o mas, la comparacion se haria sobre 6 meses. Pero enrealidad, esta solucion no resolvia el problema "de raiz", solo mitigaba el problema de que hubiera pocos meses de prueba, ademas de que la decision de que 18 meses fuese el punto de corte no tenia respaldo alguno.

Por eso fue que se corrigio y re-interpreto el enfoque de la solucion, fundamentandola en que el problema a solucionar no era el numero de meses, si no el calculo del error, es decir, que tan confiable es la comparacion de error entre los modelos. 

Fue aqui donde se propuso la prueba Diebold-Mariano, la cual consiste en una prueba estadistica que permite medir si la diferencia de error que hay entre dos modelos es significativa o es solo ruido. La prueba Diebold-Mariano requiere de un nivel de significancia, que se definio como 5\%.

Por otro lado,\texttt{MESES\_PERDIDOS\_POR\_REZAGOS} es el número fijo de meses que se pierden al calcular los rezagos (visto en la etapa anterior), es necesario para disctribuir los meses utilizables del total de meses que subió el usuario en el archivo.

Una vez definidas las variables, la primera funcion implementada en \texttt{modelos.py} es \texttt{tamano\_test\_backtest}. Esta funcion principalmente se encarga de implementar la solucion ya explicada, recibe como parametro \texttt{n\_meses\_usables}, la cantidad de meses que quedaron disponibles despues de que se hayan perdido los primeros por el calculo de los lags, y devuelve cuántos meses de esos son los que se van a usar como período de prueba (backtest).

\begin{minted}{python}
def tamano_test_backtest(n_meses_usables):
    n_meses_archivo_original = n_meses_usables + MESES_PERDIDOS_POR_REZAGOS
    if n_meses_archivo_original >= MESES_ARCHIVO_UMBRAL_TEST_GRANDE:
        return TEST_SIZE_MESES_GRANDE
    return TEST_SIZE_MESES_DEFAULT
\end{minted}

Como el umbral de desicion (\texttt{MESES\_ARCHIVO\_UMBRAL\_TEST\_GRANDE = 18}) toma en cuenta el archivo completo, y la funcion recibe la cantidad de meses disponibles ya con los 3 meses de lag descontados, con \texttt{n\_meses\_archivo\_original = n\_meses\_usables + MESES\_PERDIDOS\_POR\_REZAGOS} se le suma esos 3 meses de vuelta, para saber cuántos meses tenía el archivo que el usuario subió originalmente

Luego se compara si el total de meses completo cuenta con 18 meses, de hacelo devuelvve el tamaño de test grande con 6 meses y de no hacerlo devuelve el tamaño de test por defecto de 3 meses.

Esta función se llama desde \texttt{sistema\_prediccion.py}, para decidir el tamaño del split que se usa para elegir el campeón y generar los resultados, y desde \texttt{validacion.py}, para que los \textit{folds} de la validación cruzada temporal usen ese mismo tamaño de test y conservar un criterio estandarizado en ambas partes.

Es parte de la solucion que se propone en Thary, esta funcion no ninguna en alguna de las implementaciones que fueron revisadas ya que ninguna adapta el tamaño del perido de prueba seun la cantidad de historia disponible. 

\subparagraph{Modelos de demanda}

Las funciones que permiten entrenar y evaluar los tres modelos de demanda que se van a utilizar en el sistema: Media Móvil, Regresión Lineal y XGBoost. Estos fueron elegidos luego de realizar una revision de literatura en la que se compararon 5 modelos para la prediccion de la demanda, tal como se vio en el capitulo 2.2 "Antecedentes", de los cuales se escogieron: XgBoost y regresion lineal al ser los que presentan mejor desempeño en este tipo de problemas, y adicionalmente se incluyo media movil como un modelo mas sencillo que puede ser util para productos con demanda estable o que no cuenten con suficiente historial implementar para un modelo mas complejo.

\subparagraph{XgBoost}

La implementacion de XGBoost se hizo combinando los hiperparametros de \textit{Ejemplo 1 (SKU Demand Forecasting)} \cite{ejemplo1_demandforecast} con la transformacion del objetivo propuesta en \textit{Ejemplo 3 (XGBoost)} \cite{ejemplo3_ksmooi_inventorydemand}.

\textit{Ejemplo 3 (XGBoost)} \cite{ejemplo3_ksmooi_inventorydemand} cuenta con una configuracion muchisimo mas simple de la que necesitaria Thary, ya que en el caso del ejemplo, se estaba trabajando con aproximadamente 60 millones de filas, por lo que se priorizo la velocidad sobre el ajuste fino de hiperparamtros, por lo que no cuenta con factores como algun control de sobreajuste:

\begin{minted}{python}
# Define the parameters for the XGBoost model
params = {
    'max_depth': 8,                  # Maximum depth of a tree
    'eta': 0.5,                      # Learning rate
    'objective': 'reg:squarederror'  # Objective function for regression
}
\end{minted}

Por otra parte, el \textit{Ejemplo 1 (SKU Demand Forecasting)} \cite{ejemplo1_demandforecast} si cuenta con una configuracion mas completa que se podria reutilizar, en este se arma un XGBRegressor mas cuidado en el que se incluten variables como \texttt{learning\_rate=0.05}, \texttt{subsample=0.8}, \texttt{colsample\_bytree=0.8}, pero principalmente se escogio seguir la logica de este ejemplo porque de los cuatro ejemplos revisados, es el unico que consiste en un sistema completo con una premisa bastante similar a la de Thary, ya que ofrece pronostico de demanda, traduccion de ese pronostocp en recomendacopnes de reposicion, punto de reorden, EOQ, y todo eso bajo restricciones de lead time y presupuesto.

Es por eso, que se adopto su enfoque general completo:

\begin{itemize}
    \item Variables
    \item modelos
    \item hiperparametros
    \item comparacion de modelos
    \item logica de reposicion
\end{itemize}

Los otros ejemplos se revisaron posteriormente cuando se necesito mejorar aspectosm agregar nuevas funcionalidades o en general, agregarle mas valor y robustez a Thary.

Los hiperparámetros utilizados son los mismos que se encuentran en el \textit{Ejemplo 1 (SKU Demand Forecasting)} \cite{ejemplo1_demandforecast}. Como ya se explico anteriormente, al ser el ejemplo mas cercano a una implementacion completa de un sistema con las caracteristicas de Thary, se opto por mantener su logica de entrenamiento y ajuste de hiperparametros, con la diferencia de que \texttt{n\_estimators} se aumento a 500, \texttt{objective} se cambio a \texttt{reg:squarederror}, mientras que en el ejemplo original no se especificaba. Esto se cambio ya que se buscaba implementar el uso de minimos cuadrados para poder estimar la bondad de los modelos.

Parametros como \texttt{verbosity} fueron omitidos en Thary, pero se incluyo \texttt{n\_jobs}, el cual permite utilizar todos los nucleos de procesamiento disponibles en la maquina donde se ejecuta el entrenamiento, en lugar de un unico nucleo. Esto contribuye a reducir el tiempo de entrenamiento del modelo cuando se procesan catalogos con un numero elevado de productos.

\begin{minted}{python}
def entrenar_xgboost(X_train, y_train_log, X_test):
    model = xgb.XGBRegressor(
        objective="reg:squarederror", n_estimators=500, 
        learning_rate=0.05,
        max_depth=4, 
        subsample=0.8, 
        colsample_bytree=0.8,
        reg_alpha=0.1, 
        reg_lambda=1.0, 
        random_state=42, 
        n_jobs=-1
    )
    model.fit(X_train, y_train_log)
    pred_xgb = model.predict(X_test)
    pred_xgb = np.maximum(np.expm1(pred_xgb), 0)
    return model, pred_xgb
\end{minted}

En cuanto a la transformacion logaritmica, se implemento siguiendo la logica del \textit{Ejemplo 3 (XGBoost)} \cite{ejemplo3_ksmooi_inventorydemand}, en el que el modelo no entrena directamente sobre \texttt{y\_train}, que seria la demanda real, si no sobre \texttt{y\_train\_log}, que es la transformacion logaritmica de \texttt{y\_train}. 

Esto porque en la vida real, la demanda suele tener una distribucion en la que puede haber muchos productos con demanda baja y pocos con picos muy altos. Si se entrenara directamente sobre esos valores el modelo podria enfocarse demasiado en los pocos casos extremos, esto llevaria a que se prediciera peor la mayoria de los casos normales. El logaritmo lo que hace es parejar esa diferencia: hace que los valores muy grandes dejen de pesar tanto en comparacion con los valores mas pequeños, asi el modelo le presta atención equilibrada a todos los productos y no solo a los que tienen demanda extrema. Se usa \texttt{log1p} porque la demanda puede valer 0, y el logaritmo de 0 no existe


\subparagraph{Media Móvil.}

Media movil hace referencia al metodo mas simple de los tres. La tecnica con la que se implemento tambien se basa fuertemente en el \textit{Ejemplo 1 (SKU Demand Forecasting)} \cite{ejemplo1_demandforecast}, y consiste en simplemente reutilizar la columna de \texttt{media\_movil\_3} que ya se tienen desde la etapa anterior de features. No se tiene que hacer ningun entrenemiento. 

\begin{minted}{python}
def prediccion_baseline_movil(test_df):
    return np.maximum(test_df["media_movil_3"].values, 0)
\end{minted}

La funcion recibe como parametro el \texttt{test\_df} y posteriormente, para todas las filas, devuelve el valor \texttt{media\_movil\_3} que ya se habia calculado anteriormente.


\subparagraph{Regresion Lineal}

Para implementar regresion lineal, se adapato la logica del \textit{Ejemplo 4 (BoomBikes)} \cite{ejemplo4_boombikes}, al ser el modelo que mas pasos requiere, se implemento de manera individual en el archivo \texttt{regresion\_lineal.py}.

\begin{minted}{python}
    def _seleccionar_features_regresion(X_train, y_train, max_p_valor=0.05, max_vif=5.0):
        candidatas = [c for c in X_train.columns if X_train[c].std() > 0]
        if len(candidatas) < 2 or len(X_train) <= len(candidatas) + 1:
            return candidatas
\end{minted}

Inicialmente, se filtran y excluyen las columnas que no varian demasiado, descartando las que tienen una desviacion estandar de 0, ya que una desviacion estandar con este valor indica que en esas columnas el valor es siempre el mismo numero.

con \texttt{std() > 0}, esto se hace porque cuando se entrena regresion lineal, no siempre conviene usar todas las variables disponibles porque no todas tienen algo significativo que aportar. Es incluso perjudicial para el modelo, ya que una columna sin variacion no le sirve de nada al modelo (no hay nada que "aprender" de un número que nunca cambia).

Posteriormente, se convierten todos los valores de cada columna a un rango entre 0 y 1, para mantener las proporciones relativas con \texttt{MinMaxScaler}. Se ajusta \texttt{fit\_transform} usando solo los datos con los que se hace el entrenamiento, no se puede incluir a los de prueba porque eso haria que el sistema vea los datos que se supone debe predecir.  

\begin{minted}{python}
    try:
        scaler = MinMaxScaler()
        X_train_esc = pd.DataFrame(
            scaler.fit_transform(X_train[candidatas]),
            columns=candidatas, index=X_train.index
        )
\end{minted}

En esta parte se incluye Recursive Feature Elimination (RFE). Este consiste en un metodo que entrena varias veces un modelo de regresion lineal, y en cada una se elimina la variable que aporte menos al modelo hasta que al final, se quede con un numero controlado de variables, en este caso, tienen que ser 2 variables como minimo para posteriormente poder realizar el analisis, pero no pueden ser mas de 8 variables porque no se pueden dejar demasiadas \texttt{n\_rfe = max(2, min(8, len(candidatas) - 1))}.

Finalmente, \texttt{seleccionadas} es la lista de columnas que sobrevivieron esta primera etapa, aunque es posible que aun existan variables redundantes:

\begin{minted}{python}
    n_rfe = max(2, min(8, len(candidatas) - 1))
    rfe = RFE(LinearRegression(), n_features_to_select=n_rfe)
    rfe.fit(X_train_esc, y_train)
    seleccionadas = [c for c, usar in zip(candidatas, rfe.support_) if usar]
\end{minted}

La segunda etapa para seleccionar varibles consiste en calcular el VIF o Factor de Inflación de Varianza para cada variable. Este se encarga de medir que tan redundante es una variable en comparacion con las demas, el criterio para decidir esto es estimar si una variable se puede predecir con exactitud a partir de las otras, de ser este el caso, su VIF sera muy alto y eso es un problema porque el modelo no podria distinguir cuál de las dos variables redundantes es la que esta causando el efecto.

En la implementacion de Thary, entre las variables que se considerian redundantes se encuentran:

\begin{itemize}
    \item \texttt{inventario} y \texttt{inventario\_lag\_1}: Escencialmente tienden a ser casi la misma variable, en muchos casos, es bastante probable que el inventario de un mes sea bastante parecido al del mes anterior.
    \item Los lags (\texttt{lag\_1}, \texttt{lag\_2}, \texttt{lag\_3}, \texttt{lag\_6}): Tambien es comun que la demanda de un producto tienda a ser similar a la de meses anteriores, si un producto vendio bien un mes, es muy probable que tambien venda mes el siguiente mes, lo que a la larga hace que los lags sean variables muy parecidas entre si.
    \item (\texttt{media\_movil\_3} y \texttt{media\_movil\_6}) con los lags: Que todos los Lags sean similares llevaria a que, inevitablemente, la media movil de 3 meses y la de 6 meses tambien lo sean debido a que se calculan a partir de promediar los lags.
    \item (\texttt{media\_producto}) y las medias móviles: (\texttt{media\_producto}) consiste en el promedio historico general del producto, si los productos cuentan con un comportamiento estable, es probable que sea similar a (\texttt{media\_movil\_3} y \texttt{media\_movil\_6}), por lo que tambien se considera una variable redundante.
\end{itemize}

Teniendo esto en cuenta, las variables que no se considerarian redundantes son: \texttt{año}, \texttt{mes\_cos}, \texttt{mes\_sin}, \texttt{desviacion\_movil\_3} y \texttt{desviacion\_producto}.

\begin{minted}{python}
    while len(seleccionadas) > 1:
        X_sm = sm.add_constant(X_train_esc[seleccionadas])
        with warnings.catch_warnings():
            warnings.simplefilter("ignore", category=SingularMatrixWarning)
            warnings.simplefilter("ignore", category=UserWarning)
            vif = pd.DataFrame({
                "feature": seleccionadas,
                "VIF": [variance_inflation_factor(X_sm.values, i + 1) for i in range(len(seleccionadas))]
            })
\end{minted}

Posteriormente, se quitan las variables con un VIF infinito que nisiquiera se puede calcular, ya que de ser este el caso, significaria que la variable es tan redundante con otra que rompe las matemáticas del modelo. Esto ocurriria porque a la hora de calcular los coeficientes de la regresion lineal, el algoritmo necesita calcular cuanto aporta cada variable a la predicción, y para eso, necesita poder distinguir el efecto individual de cada una. El problema es que si dos variables son casi idénticas entre sí , es imposible distinguirlas.

Es por eso, que en esta etapa esas variables son eliminadas.En las siguientes lineas de codigo, se revisa cual variable tiene el VIF mas alto si ya no hay variables con VIF infinito, si el valor del vif mas alto supera el limite que se definio en \texttt{max\_vif}, esa variable se elimina y finalmente se vuelve a repetir el proceso desde el principio hasta que se halla calculado el VIF de las variables restantes y todas estas cuenten con un VIF menor al limite establecido.

\begin{minted}{python}
    vif_problema = vif[~np.isfinite(vif["VIF"])]
    if not vif_problema.empty:
        seleccionadas.remove(vif_problema.iloc[0]["feature"])
        continue

    peor_vif_fila = vif.loc[vif["VIF"].idxmax()]
    if peor_vif_fila["VIF"] > max_vif:
        seleccionadas.remove(peor_vif_fila["feature"])
        continue
\end{minted}

Una vez completado el proceso de eliminacion de las variables redundantes, se prosigue con la implementacion del entrenamiento del modelo de regresion lineal, en el que se revisa el "p-valor" de cada variable. 

El P-valor consiste en una medida estadistica que indica que tan probable es que la variable que se esta revisando este relacionada con la demanda solo por casualidad y no porque de verdad dicha variable este relacionada de verdad con la demanda. En este caso, si se encuentra que la variable que cuenta con el peor p-valor supera el limite de \texttt{max\_p\_valor}, que es estandarizado como 0.05 al representar el 95\% de confianza en estadistica, se elimina la variable y se repite el proceso. 

El bucle se detiene cuando ya no haya ninguna variable con el VIF alto o con el P-valor alto, al cumplir con ambos criterios, solo se contara con variables limpias y estadisticamente significativas para el proceso. 

De haber algun fallo durante esta etapa, simplemente se devuelve la lista completa de variables con el fin de garantizar de que el pipeline no se rompa y el sistema pueda continuar funcionando.

\begin{minted}{python}
        modelo_ols = sm.OLS(y_train, X_sm).fit()
        p_valores = modelo_ols.pvalues.drop("const")
        peor_p_nombre = p_valores.idxmax()
        peor_p_valor = p_valores[peor_p_nombre]

        if peor_p_valor > max_p_valor:
            seleccionadas.remove(peor_p_nombre)
        else:
            break

    return seleccionadas
except Exception:
    return candidatas
\end{minted}

Por otro lado, en ese mismo archivo, se implemento la funcion \texttt{entrenar\_regresion\_lineal}, que se encarga de entrenar el modelo de regresion lineal y devolver las predicciones, tambien es a la que se llama desde el resto del sistema.

La funcion llama a \texttt{\_seleccionar\_features\_regresion} para obtener la lista de variables que se van a usar en el entrenamiento del modelo, (si por algun motivo, dicha funcion le devuelve una lista que se encuentre vacia, se usan todas las variables originales.)

El modelo es entrenado con las variables limpias asegurandose que las predicciones no sean negativas al no ser algo que tenga sentido en este contexto. 

Finalmente, devuelve el modelo entrenado, las predicciones y la lista de variables que se usaron para entrenar el modelo.

\begin{minted}{python}
def entrenar_regresion_lineal(X_train, y_train, X_test, feature_cols):
    features_lr = _seleccionar_features_regresion(X_train, y_train)
    if not features_lr:
        features_lr = list(feature_cols)

    modelo_lr = LinearRegression()
    modelo_lr.fit(X_train[features_lr], y_train)
    pred_lr = modelo_lr.predict(X_test[features_lr])
    pred_lr = np.maximum(pred_lr, 0)

    return modelo_lr, pred_lr, features_lr

\end{minted}
